%============================================================
% Chapter 10: Mathematical Greedy
% STATUS: PLANNED  (generated by tools/outline.py; replace this file when the chapter is written)
%============================================================
\chapter{Mathematical Greedy}
\label{ch:math_greedy}

\begin{tcolorbox}[enhanced, colback=white, colframe=secondaryblue, boxrule=1pt, arc=2mm]
\textbf{\large\textcolor{primaryblue}{Greedy meets number theory: canonical coin systems, reachable sums, Egyptian fractions and greedy representations.}}
\end{tcolorbox}

\chapterfacts{1400}{2200}{3}{6}{30}

\begin{plannedbox}[Planned Contents --- this chapter is not written yet]
\textbf{Part A --- Lessons}
\begin{enumerate}[leftmargin=*, label=\textbf{A\arabic*.}, itemsep=1pt]
    \item Canonical coin systems: when greedy change-making is optimal
    \item Reachable sums: if $[1, s]$ is reachable and the next coin is $\le s+1$, then $[1, s+c]$ is
    \item Patching: the cheapest additions to make every sum reachable
    \item Greedy representations: Zeckendorf, factorial base, Egyptian fractions
    \item Greedy with divisibility and parity
\end{enumerate}
\textbf{Part B --- Patterns} (with candidate problems, checked when the chapter is written)
\begin{enumerate}[leftmargin=*, label=\textbf{P\arabic*.}, itemsep=1pt]
    \item \textbf{Reachable sums} --- Maximum Number of Consecutive Values You Can Make {\small\color{darkgray}(LeetCode 1798)}, Patching Array {\small\color{darkgray}(LeetCode 330)}
    \item \textbf{Canonical coin systems}
    \item \textbf{Greedy representations}
\end{enumerate}
\end{plannedbox}
